\documentclass[12pt,a4paper]{report}
\usepackage[bahasa]{babel}
\usepackage[utf8]{inputenc}
\usepackage[T1]{fontenc}
\usepackage{geometry}
\geometry{a4paper, margin=2.5cm, top=3cm, bottom=2.5cm, headheight=14.5pt}
\usepackage{xcolor}
\usepackage[most]{tcolorbox}
\usepackage{graphicx}
\usepackage{booktabs}
\usepackage{tabularx}
\usepackage{enumitem}
\usepackage[expansion=false]{microtype}
\usepackage{titlesec}
\usepackage{fancyhdr}
\usepackage{amsmath}
\usepackage{amssymb}
\usepackage{tikz}
\usetikzlibrary{positioning, arrows.meta}

% Warna Korporat PLN
\definecolor{plnblue}{RGB}{0,75,135}
\definecolor{plnorange}{RGB}{255,107,53}
\definecolor{lightgray}{RGB}{245,245,245}

\usepackage{hyperref}
\hypersetup{
    colorlinks=true,
    linkcolor=plnblue,
    urlcolor=plnblue,
    citecolor=plnblue
}

% Header & Footer
\pagestyle{fancy}
\fancyhf{}
\fancyhead[L]{\small\textbf{Lampiran Lembar Kerja -- Modul 4.3}}
\fancyhead[R]{\small Level 4 (Mastery) -- PLN}
\fancyfoot[C]{\small\thepage}
\renewcommand{\headrulewidth}{0.4pt}

% Format Heading
\titleformat{\chapter}[display]{\normalfont\huge\bfseries\color{plnblue}}{\chaptertitlename\ \thechapter}{20pt}{\Huge}
\titleformat{\section}{\Large\bfseries\color{plnblue}}{\thesection}{1em}{}
\titleformat{\subsection}{\large\bfseries\color{plnblue}}{\thesubsection}{1em}{}
\titlespacing*{\chapter}{0pt}{-30pt}{20pt}
\titlespacing*{\section}{0pt}{12pt}{8pt}
\titlespacing*{\subsection}{0pt}{10pt}{6pt}

% Custom Boxes
\newtcolorbox{plnbox}[1][]{%
  colback=plnblue!5!white, colframe=plnblue, fonttitle=\bfseries,
  title=#1, sharp corners, boxrule=0.8pt, breakable
}
\newtcolorbox{instruksi}{%
  colback=gray!8, colframe=gray!50!black, fonttitle=\bfseries,
  title=Instruksi Pengisian, sharp corners, breakable
}

% List styling
\setlist[itemize]{leftmargin=*, topsep=2pt, itemsep=2pt}
\setlist[enumerate]{leftmargin=*, topsep=2pt, itemsep=2pt}

% Helper: lined fill space
\newcommand{\fillline}[1][6cm]{\underline{\hspace{#1}}}
\newcommand{\writespace}[1][2cm]{\rule{0pt}{#1}}

\begin{document}

% ================= COVER =================
\begin{titlepage}
\centering
\vspace*{1.5cm}
{\Huge\bfseries\color{plnblue} LAMPIRAN}\\[0.4cm]
{\LARGE\bfseries Lembar Kerja Praktek (Worksheet)}\\[0.3cm]
{\Large\bfseries Modul 4.3: Optimalisasi Prosedur dan Alur Kerja Terintegrasi}\\[0.8cm]
\rule{0.8\linewidth}{0.6pt}\\[0.6cm]
{\large Program Penguatan Kompetensi Tata Kelola Kearsipan}\\
{\large Level 4 (Mastery)}\\[1.2cm]
{\large Disusun Khusus untuk: \textbf{PT PLN (Persero)}}\\
{\large Target Peserta: Manajemen Atas (MA)}\\[2cm]

\begin{plnbox}[Petunjuk Umum]
\begin{enumerate}[nosep, leftmargin=*]
    \item Lembar kerja ini merupakan bagian integral dari portofolio kompetensi Modul 4.3.
    \item Setiap lampiran dirancang untuk dikerjakan selama sesi workshop sesuai durasi yang telah ditetapkan.
    \item Hasil pengisian menjadi syarat penilaian kelulusan (Passing Grade $\geq$70).
    \item Kerjakan secara berkelompok (4--5 orang) kecuali dinyatakan lain.
    \item Setelah selesai, kumpulkan ke Unit Kearsipan Holding/Subholding.
\end{enumerate}
\end{plnbox}

\vfill
{\small Dicetak sebagai bagian dari paket modul pelatihan}\\
{\small\today}
\end{titlepage}

% ============================================================
%   LAMPIRAN A: Matriks Identifikasi Area Optimasi (Bab 2)
% ============================================================
\newpage
\section*{\color{plnblue} LAMPIRAN A}
\addcontentsline{toc}{section}{Lampiran A: Matriks Identifikasi Area Optimasi}
{\Large\bfseries Matriks Identifikasi Area Optimasi}\\[0.2cm]
{\normalsize\textit{Output Wajib Bab 2 -- Kriteria Penilaian 4.3.1}}\\[0.5cm]

\noindent\textbf{Nama Kelompok:} \fillline[5cm] \quad \textbf{Tanggal:} \fillline[3cm]\\[0.3cm]
\textbf{Proses yang Dievaluasi:} \fillline[9cm]\\[0.3cm]
\textbf{Unit Kerja:} \fillline[5cm] \quad \textbf{Volume Proses/Tahun:} \fillline[3cm]

\begin{instruksi}
\begin{enumerate}[nosep, leftmargin=*]
    \item Pilih 1 proses kearsipan PLN yang sedang berjalan di unit Anda.
    \item Identifikasi minimal \textbf{4 area/tahapan} yang berpotensi untuk dioptimalkan.
    \item Untuk setiap area, jelaskan jenis hambatan, potensi integrasi, dan estimasi dampak.
    \item Tentukan prioritas berdasarkan urgensi vs effort implementasi.
    \item Durasi pengerjaan: \textbf{15 menit}.
\end{enumerate}
\end{instruksi}

\vspace{0.5cm}

\begin{table}[htbp]
\centering
\small
\renewcommand{\arraystretch}{2.5}
\begin{tabularx}{\textwidth}{@{} c X X X c @{}}
\toprule
\textbf{No} & \textbf{Tahapan Proses} & \textbf{Jenis Hambatan / Titik Gesekan} & \textbf{Potensi Integrasi / Solusi} & \textbf{Prioritas} \\
\midrule
1 & & & & $\square$ T~~ $\square$ S~~ $\square$ R \\
\hline
2 & & & & $\square$ T~~ $\square$ S~~ $\square$ R \\
\hline
3 & & & & $\square$ T~~ $\square$ S~~ $\square$ R \\
\hline
4 & & & & $\square$ T~~ $\square$ S~~ $\square$ R \\
\hline
5 & & & & $\square$ T~~ $\square$ S~~ $\square$ R \\
\hline
6 & & & & $\square$ T~~ $\square$ S~~ $\square$ R \\
\bottomrule
\end{tabularx}
\end{table}

\noindent\textit{Keterangan Prioritas: T = Tinggi, S = Sedang, R = Rendah}

\vspace{0.5cm}
\noindent\textbf{Kesimpulan \& Rekomendasi Utama:}\\
\noindent\rule{\textwidth}{0.3pt}\writespace[0.8cm]\\
\noindent\rule{\textwidth}{0.3pt}\writespace[0.8cm]\\
\noindent\rule{\textwidth}{0.3pt}\writespace[0.8cm]\\
\noindent\rule{\textwidth}{0.3pt}

% ============================================================
%   LAMPIRAN B: Lembar Kerja Diagram BPMN To-Be (Bab 3)
% ============================================================
\newpage
\section*{\color{plnblue} LAMPIRAN B}
\addcontentsline{toc}{section}{Lampiran B: Lembar Kerja Diagram BPMN To-Be}
{\Large\bfseries Lembar Kerja Diagram BPMN To-Be}\\[0.2cm]
{\normalsize\textit{Output Wajib Bab 3 -- Kriteria Penilaian 4.3.2}}\\[0.5cm]

\noindent\textbf{Nama Kelompok:} \fillline[5cm] \quad \textbf{Tanggal:} \fillline[3cm]\\[0.3cm]
\textbf{Proses yang Dirancang Ulang:} \fillline[9cm]

\begin{instruksi}
\begin{enumerate}[nosep, leftmargin=*]
    \item Gunakan area kosong di bawah untuk menggambar diagram alur kerja To-Be.
    \item Gunakan notasi BPMN sederhana: $\bullet$ (Start), $\square$ (Task), $\diamond$ (Gateway), $\blacksquare$ (End).
    \item Tentukan minimal 3 swimlane (Pool/Lane) untuk peran atau sistem yang berbeda.
    \item Tandai setiap pemicu otomatis dengan warna/arsir berbeda.
    \item Durasi pengerjaan: \textbf{20 menit}.
\end{enumerate}
\end{instruksi}

\vspace{0.3cm}

% BPMN Drawing Area with swimlanes
\noindent\textbf{Diagram Alur Kerja To-Be:}
\vspace{0.3cm}

\begin{tikzpicture}
% Pool border
\draw[plnblue, thick] (0,0) rectangle (15.3cm, 14cm);
% Lane dividers
\draw[plnblue!40, dashed] (2.2cm,0) -- (2.2cm,14cm);
\draw[plnblue!30, dashed] (0,9.33cm) -- (15.3cm,9.33cm);
\draw[plnblue!30, dashed] (0,4.67cm) -- (15.3cm,4.67cm);
% Lane labels
\node[rotate=90, font=\small\bfseries, text=plnblue] at (1.1cm, 11.67cm) {Lane 1:};
\node[rotate=90, font=\small\itshape, text=gray] at (1.1cm, 10.3cm) {\fillline[2cm]};
\node[rotate=90, font=\small\bfseries, text=plnblue] at (1.1cm, 7cm) {Lane 2:};
\node[rotate=90, font=\small\itshape, text=gray] at (1.1cm, 5.6cm) {\fillline[2cm]};
\node[rotate=90, font=\small\bfseries, text=plnblue] at (1.1cm, 2.33cm) {Lane 3:};
\node[rotate=90, font=\small\itshape, text=gray] at (1.1cm, 0.9cm) {\fillline[2cm]};
\end{tikzpicture}

\vspace{0.3cm}
\noindent\textbf{Daftar Pemicu Otomatis:}

\begin{table}[htbp]
\centering
\small
\renewcommand{\arraystretch}{2.0}
\begin{tabularx}{\textwidth}{@{} c X X X @{}}
\toprule
\textbf{No} & \textbf{Pemicu / Status} & \textbf{Aksi Otomatis} & \textbf{Sistem Sumber} \\
\midrule
1 & & & \\
\hline
2 & & & \\
\hline
3 & & & \\
\bottomrule
\end{tabularx}
\end{table}

% ============================================================
%   LAMPIRAN C: Template Draft SOP Digital (Bab 4)
% ============================================================
\newpage
\section*{\color{plnblue} LAMPIRAN C}
\addcontentsline{toc}{section}{Lampiran C: Template Draft SOP Digital Terintegrasi}
{\Large\bfseries Template Draft SOP Digital Terintegrasi}\\[0.2cm]
{\normalsize\textit{Output Wajib Bab 4 -- Kriteria Penilaian 4.3.3}}\\[0.5cm]

\noindent\textbf{Nama Kelompok:} \fillline \quad \textbf{Tanggal:} \fillline[3cm]

\begin{instruksi}
\begin{enumerate}[nosep, leftmargin=*]
    \item Isi ketujuh komponen SOP untuk 1 proses kearsipan PLN yang telah dirancang di Bab 3.
    \item Pastikan setiap komponen merujuk pada regulasi dan standar yang relevan.
    \item Lakukan peer review silang dengan kelompok lain setelah selesai.
    \item Durasi pengerjaan: \textbf{15 menit}.
\end{enumerate}
\end{instruksi}

\vspace{0.4cm}

% Komponen 1
\begin{tcolorbox}[colback=plnblue!3, colframe=plnblue!50, sharp corners, title={\small\bfseries Komponen 1: Tujuan \& Ruang Lingkup}]
\vspace{2.2cm}
\end{tcolorbox}

% Komponen 2
\begin{tcolorbox}[colback=plnblue!3, colframe=plnblue!50, sharp corners, title={\small\bfseries Komponen 2: Definisi \& Akronim}]
\vspace{2.2cm}
\end{tcolorbox}

% Komponen 3
\begin{tcolorbox}[colback=plnblue!3, colframe=plnblue!50, sharp corners, title={\small\bfseries Komponen 3: Peran \& Tanggung Jawab}]
\small
\renewcommand{\arraystretch}{2.0}
\begin{tabularx}{\textwidth}{@{} X X X @{}}
\toprule
\textbf{Peran} & \textbf{Jabatan/Fungsi} & \textbf{Tanggung Jawab Utama} \\
\midrule
Inisiator & & \\
\hline
Validator & & \\
\hline
Custodian & & \\
\hline
Approver & & \\
\bottomrule
\end{tabularx}
\renewcommand{\arraystretch}{1.0}
\end{tcolorbox}

% Komponen 4
\begin{tcolorbox}[colback=plnblue!3, colframe=plnblue!50, sharp corners, title={\small\bfseries Komponen 4: Pemicu Sistem \& Metadata Wajib}]
\noindent\textbf{Pemicu:} \fillline[10cm]\\[0.3cm]
\textbf{Sistem Sumber:} \fillline[10cm]\\[0.3cm]
\textbf{Metadata Wajib (min. 8 field):}\\[0.2cm]
\begin{tabularx}{\textwidth}{@{} X X X X @{}}
1. \fillline[3cm] & 2. \fillline[3cm] & 3. \fillline[3cm] & 4. \fillline[3cm] \\[0.4cm]
5. \fillline[3cm] & 6. \fillline[3cm] & 7. \fillline[3cm] & 8. \fillline[3cm] \\
\end{tabularx}
\end{tcolorbox}

\newpage

% Komponen 5
\begin{tcolorbox}[colback=plnblue!3, colframe=plnblue!50, sharp corners, title={\small\bfseries Komponen 5: Alur Validasi \& SLA}]
\small
\renewcommand{\arraystretch}{2.0}
\begin{tabularx}{\textwidth}{@{} X X c X @{}}
\toprule
\textbf{Tahapan} & \textbf{PIC} & \textbf{SLA} & \textbf{Eskalasi Jika Lewat} \\
\midrule
Validasi Level 1 & & jam & \\
\hline
Validasi Level 2 & & jam & \\
\hline
Approval Final & & jam & \\
\bottomrule
\end{tabularx}
\renewcommand{\arraystretch}{1.0}
\end{tcolorbox}

% Komponen 6
\begin{tcolorbox}[colback=plnorange!5, colframe=plnorange!50, sharp corners, title={\small\bfseries Komponen 6: Exception Handling \& Fallback}]
\noindent\textbf{Skenario 1 -- Sistem Down:}\\
\rule{\textwidth}{0.3pt}\writespace[0.7cm]\\
\rule{\textwidth}{0.3pt}\writespace[0.5cm]

\noindent\textbf{Skenario 2 -- Validasi Ditolak:}\\
\rule{\textwidth}{0.3pt}\writespace[0.7cm]\\
\rule{\textwidth}{0.3pt}\writespace[0.5cm]

\noindent\textbf{Skenario 3 -- Metadata Tidak Lengkap:}\\
\rule{\textwidth}{0.3pt}\writespace[0.7cm]\\
\rule{\textwidth}{0.3pt}
\end{tcolorbox}

% Komponen 7
\begin{tcolorbox}[colback=plnblue!3, colframe=plnblue!50, sharp corners, title={\small\bfseries Komponen 7: Referensi Regulasi \& Standar}]
\begin{enumerate}[nosep, leftmargin=*]
    \item \fillline[12cm]
    \item \fillline[12cm]
    \item \fillline[12cm]
    \item \fillline[12cm]
\end{enumerate}
\end{tcolorbox}

\vspace{0.5cm}
\noindent\textbf{Catatan Peer Review (diisi oleh kelompok penilai):}\\
\noindent\rule{\textwidth}{0.3pt}\writespace[0.8cm]\\
\noindent\rule{\textwidth}{0.3pt}\writespace[0.8cm]\\
\noindent\rule{\textwidth}{0.3pt}

% ============================================================
%   LAMPIRAN D: Business Case & Action Plan 30 Hari (Bab 5)
% ============================================================
\newpage
\section*{\color{plnblue} LAMPIRAN D}
\addcontentsline{toc}{section}{Lampiran D: Business Case \& Action Plan 30 Hari}
{\Large\bfseries Business Case \& Action Plan 30 Hari}\\[0.2cm]
{\normalsize\textit{Output Wajib Bab 5 -- Kriteria Penilaian 4.3.4}}\\[0.5cm]

\noindent\textbf{Nama Kelompok:} \fillline \quad \textbf{Tanggal:} \fillline[3cm]\\[0.3cm]
\textbf{Proses yang Dioptimalkan:} \fillline[10cm]\\[0.3cm]
\textbf{Unit Kerja Target Pilot:} \fillline[10cm]

\begin{instruksi}
\begin{enumerate}[nosep, leftmargin=*]
    \item Isi seluruh komponen business case berdasarkan analisis di Bab 3--5.
    \item Gunakan data kalkulasi ROI yang telah dihitung untuk bagian estimasi manfaat.
    \item Susun action plan 30 hari yang realistis dengan PIC jelas.
    \item Durasi pengerjaan: \textbf{15 menit}.
\end{enumerate}
\end{instruksi}

\vspace{0.3cm}

% Business Case
\begin{tcolorbox}[colback=plnblue!3, colframe=plnblue, sharp corners, title={\bfseries Business Case 1 Halaman}]

\noindent\textbf{1. Latar Belakang \& Masalah:}\\
\rule{\textwidth}{0.3pt}\writespace[0.7cm]\\
\rule{\textwidth}{0.3pt}\writespace[0.7cm]\\
\rule{\textwidth}{0.3pt}\writespace[0.5cm]

\noindent\textbf{2. Solusi yang Diusulkan:}\\
\rule{\textwidth}{0.3pt}\writespace[0.7cm]\\
\rule{\textwidth}{0.3pt}\writespace[0.5cm]

\noindent\textbf{3. Target KPI:}

\vspace{0.2cm}
{\small
\renewcommand{\arraystretch}{1.8}
\begin{tabularx}{\textwidth}{@{} X c c @{}}
\toprule
\textbf{Indikator} & \textbf{Baseline (Saat Ini)} & \textbf{Target (30 Hari)} \\
\midrule
Waktu Siklus & & \\
\hline
Ketepatan Awal & & \\
\hline
Tingkat Kepatuhan Prosedur & & \\
\hline
Waktu Temu Kembali & & \\
\bottomrule
\end{tabularx}
\renewcommand{\arraystretch}{1.0}
}

\noindent\textbf{4. Estimasi Manfaat Finansial:}\\[0.2cm]
\begin{tabularx}{\textwidth}{@{} X r @{}}
Penghematan waktu siklus: & Rp \fillline[4cm] /tahun \\[0.3cm]
Reduksi biaya error/revisi: & Rp \fillline[4cm] /tahun \\[0.3cm]
Mitigasi risiko kepatuhan: & Rp \fillline[4cm] /tahun \\[0.3cm]
\textbf{Total Estimasi Manfaat:} & \textbf{Rp \fillline[4cm] /tahun} \\
\end{tabularx}

\vspace{0.3cm}
\noindent\textbf{5. Estimasi Biaya Implementasi:} Rp \fillline[4cm]\\[0.3cm]
\noindent\textbf{6. ROI Proyeksi:} \fillline[3cm]\% \quad \textbf{Payback Period:} \fillline[3cm]

\end{tcolorbox}

\vspace{0.4cm}

% Action Plan 30 Hari
\begin{tcolorbox}[colback=plnorange!5, colframe=plnorange, sharp corners, title={\bfseries Action Plan 30 Hari}]

\small
\renewcommand{\arraystretch}{2.2}
\begin{tabularx}{\textwidth}{@{} c X X c @{}}
\toprule
\textbf{Minggu} & \textbf{Aktivitas Utama} & \textbf{PIC} & \textbf{Status} \\
\midrule
\textbf{1} & & & $\square$ \\
\hline
\textbf{2} & & & $\square$ \\
\hline
\textbf{3} & & & $\square$ \\
\hline
\textbf{4} & & & $\square$ \\
\bottomrule
\end{tabularx}
\renewcommand{\arraystretch}{1.0}

\noindent\textbf{Mekanisme Monitoring:} \fillline[8cm]\\[0.3cm]
\noindent\textbf{Pelaporan Hasil Kepada:} \fillline[8cm]\\[0.3cm]
\noindent\textbf{Tanggal Laporan Akhir:} \fillline[8cm]

\end{tcolorbox}

\vspace{0.5cm}

\begin{center}
\begin{tcolorbox}[colback=green!5, colframe=green!50!black, width=0.85\textwidth, sharp corners]
\centering
\textbf{Tanda Tangan Kelompok}\\[0.5cm]
\begin{tabularx}{0.9\textwidth}{@{} X X X @{}}
\centering Ketua Kelompok & \centering Anggota 1 & \centering Anggota 2 \tabularnewline[1.5cm]
\centering\rule{3cm}{0.3pt} & \centering\rule{3cm}{0.3pt} & \centering\rule{3cm}{0.3pt} \tabularnewline[0.2cm]
\centering Anggota 3 & \centering Anggota 4 & \centering Fasilitator \tabularnewline[1.5cm]
\centering\rule{3cm}{0.3pt} & \centering\rule{3cm}{0.3pt} & \centering\rule{3cm}{0.3pt} \tabularnewline
\end{tabularx}
\end{tcolorbox}
\end{center}

\end{document}
